The HQC.KEM is obtained starting from the HQC.PKE, augmenting it with components
which do not involve polynomial arithmetic operations. 
In particular, the HQC.KEM employs the \textsc{SHAKE} extendable output function~\cite{SHAandSHAKE} (i.e., a function mapping an arbitrary-length bit 
string to a bitstring with a user chosen length, which is indistinguishable from
a randomly selected string; in other words a PRNG algorithm with a seed of 
arbitrary length) to realize the variant of the FO 
transform~\cite{DBLP:conf/tcc/HofheinzHK17}. Since the constant time and 
side channel protection of \textsc{SHAKE} have been successfully tackled in the 
literature, we consider them out of scope for this work, 
and employ the existing results.
Except for \textsc{SHAKE}, the HQC.KEM algorithmic components are the same of the 
HQC.PKE, namely polynomial addition and multiplication in $\mathcal{R}$, 
pseudorandom bit generation via \textsc{SHAKE} and encoding/decoding for the 
concatenated Reed-M\"{u}ller Reed-Solomon code.
Considering the fact that binary polynomial additions are straightforward to 
compute in constant time, and protect with respect to side channel attacks, 
we focus on polynomial multiplications and encoding/decoding. The pseudorandom
number generation component is effectively and efficiently protected in the
same way the HQC.KEM \textsc{SHAKE} element is.

Concerning the encoding/decoding, we devise an algorithmic strategy to avoid
the requirement for a power based side channel attack protected decoder, 
exploiting the linearity of the RM-RS codes, while no timing analysis concerns 
are present thanks to their constant time implementation being the current one 
at the state of the art.
In the following, we describe in detail the approaches to a timing-side channel
protected polynomial multiplication, while deferring the protection of the said
multiplication against power analysis attacks to the description of our 
algorithmic strategy for the protection of the entire HQC.PKE.

\subsection{Polynomial Multiplication Approaches}
HQC.PKE uses polynomial multiplications in $\mathcal{R}=\mathbb{F}_2[x] / (x^n-1)$
in each of its primitives. Referring to the definition of the HQC.PKE 
primitives and to Table~\ref{tab:params}, \textsf{Keygen()} requires one 
multiplication between elements in $\mathcal{R}$, where the polynomial 
employed as a second factor has only $\omega$ of non-null coefficients, 
with $\omega \ll n$, i.e., $\mathbf{h} \cdot \mathbf{y}$, with 
$\mathbf{h} \in \mathcal{R}$, $\mathbf{y} \in \mathcal{R}_{\omega}$. 
Thus, in any software implementation, the first factor (e.g., $\mathbf{h}$) 
is stored as a {\em dense} polynomial, i.e., a sequence of $n$ binary 
coefficients, while the second one ($\mathbf{y}$) is stored as a {\em sparse} 
polynomial, i.e., as a sequence of integer numbers, each of which is intended 
as the exponent of a non-null monomial ($x^i, 0 \leq i \leq n-1$). 
\textsf{Encrypt()} requires two multiplications in $\mathcal{R}$ to compute the 
ciphertext, in each of which the first factor is a dense polynomial and the 
second factor is a sparse polynomial with a number of coefficients equal to 
$\omega_r$. 
Finally, \textsf{Decrypt()} requires one multiplication in $\mathcal{R}$, where 
the first factor is a dense polynomial ($\mathbf{u}$) and the second one is a 
sparse polynomial ($\mathbf{y}$) with $\omega$ non-null coefficients.\\ 
Denoting as $\mathtt{wsize}$ the bitlength of the word of the CPU employed as a 
target computational platform, the data structure for storing a dense polynomial 
$b(x) \in \mathcal{R}$ is an array $B = (B_{t-1}, \ldots, B_j, \ldots, B_0)_{\mathtt{w}}$ of $t = \left\lceil \frac{n}{\mathtt{wsize}} \right\rceil$ cells, where each cell 
$B_j$ stores a sequence of adjacent binary coefficients in $b(x)$ that fits in 
a CPU word. Equivalently, the same polynomial can be thought of as a 
multiple precision unsigned integer encoded in $t$ CPU words, with each digit 
(limb) $0$$\leq$$B_j$$<$$\mathtt{w}$, $0$$\leq$$j$$\leq$$t$$-$$1$, meant to 
follow the usual positional notation of a number in radix 
$\mathtt{w} = 2^{\mathtt{wsize}}$.
A sparse polynomial $a(x) \in \mathcal{R}$ is an array $A = (A_{\omega-1}, \ldots, A_i, \ldots, A_{0})$, with $\omega \ll n$ cells, such that each cell $A_i=a$ is an 
unsigned integer that encodes the exponent of a non-null monomial 
$x^a \in \mathcal{R}$ included in $a(x)$, where the value $0\leq a\leq n-1$ is 
encoded with at most $m=\lceil \log_2(n)\rceil$ bits. 
The product $c(x) = b(x) \cdot x^{A_i} \bmod (x^n-1)$ between a dense polynomial 
$b(x)$ and a monomial with $A_i = a$  must be equivalent to the computation of 
a circular left-shift of the coefficients of $b(x)$, for $a$ positions across 
their original $n$-bit positions (also denoted as $\clshift a$). 

\subsection{Sparse-to-dense Multiplication Algorithm}\label{subsec:shift_add}
% \begin{figure*}[!t]
% \noindent\begin{minipage}{0.99\textwidth}
% \footnotesize
% \begin{algorithm}[H]
%     \caption{Sparse-to-dense Polynomial Multiplication\label{alg:varTimeShiftAndAdd}}
%     \KwIn{$n$: maximum number of bits for a polynomial in $\mathcal{R}$$=$$\mathbb{F}_2[x]/(x^n$$-$$1)$.\newline 
%     $\mathtt{wsize}$: bitlength of a CPU word, $\mathtt{w} = 2^{\mathtt{wsize}}$ radix employed to represent in positional notation a long integer/binary-polynomial.\newline   
%     $A = (A_{\omega-1}, \ldots, A_i, \ldots, A_{0})$, represents $a(x) \in \mathcal{R}$, where $A_i$ is an unsigned integer that encodes the exponent of a non-null monomial $x^a \in \mathcal{R}$, $0\leq a\leq n-1$, $\omega < n$.\newline 
%           $B = (B_{t-1}, \ldots, B_j, \ldots, B_0)_{\mathtt{w}}$, $t = \left\lceil \frac{n}{\mathtt{wsize}} \right\rceil$, represents $b(x) \in \mathcal{R}$, where a limb $B_j$ stores a sequence of binary coefficients fitting in a CPU word having $\mathtt{wsize}$ bits (i.e., $0$$\leq$$B_j$$<$$\mathtt{w}$), $0$$\leq$$j$$\leq$$t$$-$$1$.} 
%     \KwOut{$C = (C_{t-1}, \ldots, C_k, \ldots, C_0)_{\mathtt{w}}$, represents $c(x) = a(x)\cdot b(x) \in \mathcal{R}$, where each limb $C_k$ stores a sequence of binary coefficients fitting into a CPU word (i.e., $0$$\leq$$C_k$$<$$\mathtt{w}$), $0$$\leq$$k$$\leq$$t$$-$$1$.}
%     \BlankLine
%     $C \gets (0,\ldots,0)$ \tcp*[l]{\color{blue}zeroed array with $t$ cells, each with\ \text{{\em wsize}}\ bits}
%     \For{$i \gets 0$ \KwTo $\omega-1$}{
%        $\mathtt{tmp} \gets (\ B \clshift A_i\ )$\tcp*[l]{\color{blue}\footnotesize variable-time comput. of a $t$-word array}    
%        \For{$k \gets 0$ \KwTo $(t-1)$} {
%            $C_k \gets C_k \oplus \mathtt{tmp}_k$\;
%         }
%     }
%     \BlankLine
%     \KwRet $C$\;
%     \end{algorithm}
% \end{minipage}
% \vspace{-0.4cm}
% \end{figure*}
In the third round of the NIST standardization process, the authors of HQC 
proposed, as a reference implementation, only a schoolbook non-constant time 
polynomial multiplication algorithm using a minor variant of the schoolbook
\textit{Shift-and-Add} method, exploiting the sparsity of one of the two 
operands.
The approach is summarized as $c(x) = b(x) \cdot a(x) = b(x) (x^{A_{\omega-1}}\oplus\cdots \oplus x^{A_{i}}\oplus\ldots)$ $= b(x)x^{A_{\omega-1}} \oplus \cdots\oplus b(x)x^{A_{i}} \cdots\oplus b(x)x^{A_{0}}$, where $A_{\omega-1},\ldots, A_{0} \in \{0,\ldots ,n-1\}$ 
and each multiplication by a monomial is computed as a circular left-shift of 
the coefficients of the first polynomial $b(x)$ for $A_i$ positions across its 
original $n$ bit positions. 
The pseudo-code of the method is in Algorithm~\ref{alg:varTimeShiftAndAdd}, Appendix~\ref{sec:appendix}, while its worst-case computational 
complexity is ${O}(\omega(2t)$) bit operations.

The major drawback of the schoolbook \textit{Shift-and-Add} method lies in the 
variable-time computation reported on Alg.~\ref{alg:varTimeShiftAndAdd}, 
$\mathbf{line}\ 3$, which reveals on the timing side-channel, the sensitive 
value of each monomial coefficient $A_i$ of the sparse polynomial factor. 
Indeed, the computation time of a circular left-shift for $a$ positions across 
the original $n$ bit positions of the first array is a multiple of the execution 
time required to compute ($B \clshift 1$).         
% 
\begin{figure*}[!t]
\noindent\begin{minipage}{0.99\textwidth}
\footnotesize
\begin{algorithm}[H]
\SetAlgoLined
    \caption{Constant-time Sparse-to-dense Polynom. Multiplication\label{alg:CTShiftAndAdd}}
    \KwIn{$n$: maximum number of bits for a polynomial in $\mathcal{R}$$=$$\mathbb{F}_2[x]/(x^n$$-$$1)$.\newline 
    $\mathtt{wsize}$: bitlength of a CPU word, $\mathtt{w} = 2^{\mathtt{wsize}}$ radix employed to represent in positional notation a long integer or the sequence of coefficients of a binary-polynomial.\newline   
    $A = (A_{\omega-1}, \ldots, A_i, \ldots, A_{0})$, represents $a(x) \in \mathcal{R}$, where $A_i=a$ is an unsigned integer that encodes the exponent of a non-null monomial $x^a \in \mathcal{R}$, $0\leq a\leq n-1$, $\omega < n$, $a = (a_{m-1}, \ldots, a_0)_{\mathtt{bin}}$, $a_i \in \{0,1\}$, $m = \lceil \log_2(n) \rceil$.\newline 
    $B = (B_{t-1}, \ldots, B_j, \ldots, B_0)_{\mathtt{w}}$, $t = \left\lceil \frac{n}{\mathtt{wsize}} \right\rceil$, represents $b(x) \in \mathcal{R}$, where a limb $B_j$ stores a sequence of binary coefficients fitting in a CPU word having $\mathtt{wsize}$ bits (i.e., $0$$\leq$$B_j$$<$$\mathtt{w}$), $0$$\leq$$j$$\leq$$t$$-$$1$.} 
    \KwOut{$C = (C_{t-1}, \ldots, C_k, \ldots, C_0)_{\mathtt{w}}$, represents $c(x) = x^a \cdot b(x) \in \mathcal{R}$, where each limb $C_k$ stores a sequence of binary coefficients fitting into a CPU word (i.e., $0$$\leq$$C_k$$<$$\mathtt{w}$), $0$$\leq$$k$$\leq$$t$$-$$1$.}
    \BlankLine
    \For{$i \gets 0$ \KwTo $\omega-1$}{
    $a \gets A_i$\tcp*[l]{\color{blue}value with $m$$=$$\lceil \log_2(n)\rceil$ bits: $(a_{m-1},\ldots,m_0)_{\mathtt{bin}}$}
    $C \gets (\ B \clshift (a \mod \mathtt{wsize}) \ )$\tcp*[l]{\color{blue}\footnotesize computation of a $t$-word array}
    \For{$j \gets \log_2(\mathtt{wsize})$ \KwTo $(m-1)$} {
        $\mathtt{tmp} \gets (\ B \clshift 2^j\ )$\tcp*[l]{\color{blue}\footnotesize computation of a $t$-word array}
        $\mathtt{mask} \gets 0 - (a_j = 1)$\tcp*[l]{\color{blue}\footnotesize single word: all its bits are $1$ or $0$}
        \For{$k \gets 0$ \KwTo $(t-1)$} {
           $C_k \gets C_k \oplus ((C_k \wedge \mathtt{mask}) \vee (\mathtt{tmp}_k \wedge \mathtt{mask}))$\;
        }
    }
    }
    \KwRet $C$\;
    \end{algorithm}
\end{minipage}
\vspace{-0.4cm}
\end{figure*}
A constant-time realization of the Shift-and-Add strategy is shown in Algorithm~\ref{alg:CTShiftAndAdd}. The rationale is still to split the computation 
as the iterated sum of the dense polynomial multiplied by a monomial one. 
However, instead of executing in variable-time every single monomial 
multiplication, it is possible to realize it as a fixed sequence of circular 
left-shift operations followed by arithmetic predicates for the conditional 
committment of each of them.
Specifically, considering $a$$=$$(a_{m-1}, \ldots, a_j, \ldots, a_0)_{\mathtt{bin}}$, $0$$\leq$$j$$<$$m$, $a_j$$\in$$\{0,1\}$, $m$$=$$\lceil \log_2(n) \rceil$, 
the multiplication $c(x) = b(x) \cdot x^a \bmod (x^n$$-$$1)$ is decomposed as: 
\begin{center}
$c(x) = b(x)  \cdot \left( x^{a_{m-1}2^{m-1}} \cdots x^{a_{j}2^{m-1}} \cdots x^{a_{0}} \right)\bmod (x^n-1) $ $=\phantom{MMMMMMM...}$\newline
$= b(x)  \cdot \left( x^{a_{m-1}2^{m-1}} \cdots x^{a_{\log_2(\mathtt{wsize})}2^{\log_2(\mathtt{wsize})}} \cdots x^{a \bmod \mathtt{wsize}} \right) \bmod (x^n-1)$, 
\end{center}
observing that $\mathtt{wsize}$ is a power of two (for any CPU), thus $a = (a_{m-1},  \ldots, a_0)_{\mathtt{bin}}$ $= (a_{m-1}, \ldots, a_{\log_2(\mathtt{wsize})})_{\mathtt{bin}}$ concatenated with the binary string encoding $(a \bmod \mathtt{wsize})$, while the execution of $b(x) \cdot x^{a \bmod \mathtt{wsize}}$ is amenable to be realized through a constant amount of $t=\lceil \frac{n}{\mathtt{wsize}}\rceil$ circular left-shifts. These observations, together with  the arithmetic test to establish whether any $a_j$ is asserted or not, and a conditional update of the final result (see Alg.~\ref{alg:CTShiftAndAdd}: \textbf{lines} $4$--$8$), allow to realize a constant-time version of the Shift-and-Add method with a computational cost that amounts to $\Theta\left(\omega\,(2t)\,m + \omega\, t+\omega\, m\right)$ bit operations. 

\subsection{Dense-to-dense Multiplication Algorithms}\label{subsec:CTDenseMul}
Another approach to constant-time polynomial multiplication is to give up the 
savings offered by the sparse representation of one of the factors and represent
both factors as dense polynomials. 
Indeed, the execution of algorithmic strategies as the schoolbook multiplication 
(recommended for operands encoded with up to $\approx 300$ bits), the 
Karatsuba-, Toom-Cook- or Fast-Fourier-Transform-based (FFT) multiplications 
(recommended for operands with bitlength up to $\approx 10^4$, from 
$\approx 10^4$ to $\approx 6\cdot10^4$, and $\gtrapprox 6\cdot10^4$, 
respectively) exhibit no information leakage on the timing side-channel. 
 
\noindent
\textbf{HQC Approach.}
% From a performance perspective, the choice of the algorithmic strategy that better fits the multiplication operations listed by the HQC primitives depends on the bitlength of the operands, which in turn depend on the HQC instance that is specifically considered among \texttt{hqc-128}, \texttt{hqc-192}, and \texttt{hqc-256}.
In the fourth round of the NIST standardization process, given the size of the 
polynomials listed in Table~\ref{tab:params}, for the realization of polynomial 
multiplications required in \texttt{hqc-128}, \texttt{hqc-192}, the authors of 
HQC decided to apply a combination of the Toom-Cook multiplication strategy 
named as Toom-$3$, which splits each factor in three parts before performing 
operations on them, with a recursive $3$-split Karatsuba multiplication approach 
or a recursive $5$-split Karatsuba multiplication approach, down to the point 
where a $64$-bit$ \times 64$-bit multiplier is sufficient. 
For the realization of polynomial multiplications required in \texttt{hqc-256} 
the same overall approach is applied, replacing the $3$-split Karatsuba 
algorithm with a $5$-split one. The chosen strategy is supported by a favorable 
$O(n^{1.46})$ cost of the Toom-Cook approach, which is further optimized, in 
the finite regime of the HQC parameters, applying also the Karatsuba technique 
and the (vectorized) multiply instructions provided by the instruction set 
architecture of the target platform.
 
\noindent
\textbf{Alternate Approach.}
The asymptotically fastest methods currently known to perform multiplications in 
$\mathcal{R}$ require to evaluate them at $2n$ points, perform a pointwise 
multiplication of the these results, and interpolate the resulting $2n$ values 
to get the final polynomial. While this approach may seem less efficient, 
with a proper choice of the evaluation points it takes the name of 
Fast Fourier Transform (FFT) based multiplication, and the overall 
computational cost (including the one of interpolation) requires 
${O}(n\cdot\log(n))$ bit operations.

We noted that relatively recent development in the field of FFT-based arithmetic 
in binary finite fields may allow execution of polynomial multiplications in 
both $\mathtt{F}_2[x]$ and $\mathcal{R}$ exhibiting a computational complexity 
that is less than the one that can be obtained with other methods when the 
bitlength of the operands in in the range prescribed by each HQC instance.
% 
Specifically, Cantor~\cite{DBLP:journals/jct/Cantor89}, developed the so-called 
additive-FFTs through evaluating polynomials at points that form an additive 
subgroup in $\mathbb{F}_2$ and defining a basis, termed {\em Cantor basis}, for 
constructing an FFT over the field that works with the complexity of 
$O(n \log(n))$ multiplications and 
$O(n\cdot\log^{\log_2 (3)}(n))$ $\mathtt{xor}$-additions for evaluating 
$n = 2^m$ elements. 
In $2010$, Gao and Mateer~\cite{DBLP:journals/tit/GaoM10} presented an 
additive-FFT (GM-FFT) over $\mathbb{F}_{2^k}$ with 
$O(n\cdot\log^2(n))$ $\mathtt{xor}$-operations for evaluating a polynomial at 
$n = 2^m $ points in general. 
It can be optimized to $O(n\cdot\log( n)\cdot\log (\log(n)))$ only when $m$ is 
a power of $2$, and the polynomials in this case are represented in a special 
basis introduced by Cantor~\cite{DBLP:journals/jct/Cantor89}. 
In $2014$, Bernstein and Chou~\cite{DBLP:conf/sacrypt/BernsteinC14} presented an 
efficient implementation of the GM-FFT with performance equal to the ones 
obtainable applying the more general variant introduced in the same year by 
Lin, Chung, and Han~\cite{DBLP:conf/focs/LinCH14} (named LCH-FFT ), which uses 
a polynomial basis different from the standard one. 
In the subsequent work~\cite{DBLP:journals/tit/LinAH16} a method to convert 
from/to the standard polynomial basis to/from the one required 
by~\cite{DBLP:journals/jct/Cantor89} is introduced. 
The computational complexity to realize the conversion is 
$O(n\cdot \log(n)\cdot \log(\log(n)))$ $\mathtt{xor}$-operations, with $n$ 
being any power of two, which in turn reaches the same complexity as the 
common multiplicative FFTs.

We were able to effectively reproduce the multiplication strategy introduced  
in~\cite{DBLP:conf/sacrypt/BernsteinC14} to the end of realizing the polynomial 
multiplications listed in the HQC primitives with a constant-time dense-to-dense
algorithmic approach.
Since the length of the operands $n$ in HQC is not a power of two, a 
straightforward approach would call for zero-padding of the polynomials 
for a number of coefficients up to the nearest power of two greater than $n$.
However, this approach performs a significant amount of unnecessary computations
among zero-filled words, whenever $n$ is far from the said power of two (as it
is the case in \textsc{hqc-128} and \textsc{hqc-192}).
To avoid such overhead, we automatically generated the code of the  
multiplication algorithm in~\cite{DBLP:conf/sacrypt/BernsteinC14}, removing at code
generation time the operations acting deterministically only on zero-filled words.\footnote{The source code of our HQC realization will be publicly available upon 
acceptance} 

\subsection{Boolean Masking of the HQC.PKE against Power Analysis}\label{subsec:masking}
Regardless of their security against mathematical cryptanalysis, the 
implementation of cryptographic algorithms on end-user devices may leak 
information about secret parameters of the scheme, which can lead to 
successful plaintext or secret-key recovery attacks. 
Gathering power- or EM-measurements (a.k.a. {\em traces}) during a known-input 
execution of a cryptographic primitive allows to apply a divide-and-conquer 
strategy in inferring the secret values in the computation, thanks to the 
timewise locality and granularity of the information obtained.
The general workflow for non-profiled side channel attacks predicts the side
channel behaviour of a simple operation in the cryptographic primitive, depending
on a (guessed) value for the small portion of secret involved in it.
After collecting measurements from the device, the correct guess for the secret
is singled out comparing with statistical means the secret-dependent predictions
with the actual measures.

A popular, and quantifiably secure countermeasure against side-channel attacks 
is \emph{masking}~\cite{ishai2003private,prouff2013masking}.
Masking works through performing the entire computation on share-split values, 
through an $x$-out-of-$x$ share splitting schemes.
The number of shares into which a value is split provides a quantification
of the security margin: this is commonly denoted as the \emph{masking order} $d$, 
where a masking scheme with $x$ shares is known as a $d=x-1$ order masking.
The most common share splitting strategy is the one known as 
\emph{Boolean masking}.
Applying Boolean masking to a polynomial in $\mathbf{v}\in\mathcal{R}$ we have that the $x$-uple $\ssplit{\mathbf{v}}= (\mathbf{v}_0,\ldots,\mathbf{v}_{x-1})$ of the shares of a polynomial $\mathbf{v}\in \mathcal{R}$ is obtained drawing $x-1$ random polynomials, 
$(\mathtt{rnd}_0,\ldots,\mathtt{rnd}_{x-2})$ and obtaining the remaining share 
as their coefficient-wise \texttt{xor} sum with the polynomial to be 
split $\mathbf{v}$, $(\mathbf{v}\oplus \bigoplus_{i=0}^{x-2} 
\mathtt{rnd}_i)$. We denote this operation as 
$\text{\sc Split}(\mathbf{v}) = (\mathbf{v}\oplus \sum_{i=0}^{x-2} \mathtt{rnd}_i,\mathtt{rnd}_0,\ldots,\mathtt{rnd}_{x-2})=\ssplit{\mathbf{v}}$.

All the operations on share-split polynomials must be performed so that combining
the share-split results (i.e., \texttt{xor}-ing together all shares) yields the correct result. 
We will denote such an action on the share-split value $\ssplit{\mathbf{v}}$ as
$\text{\sc Combine}(\ssplit{\mathbf{v}})=\mathbf{v}$.
The share combination is performed only at the end of the primitive computation, as any operation combining all the shares from of a single original value $\mathbf{v}$ reveals information on the value itself.
By contrast, combining up to $x-1$ shares does not yield any information 
leakage.

Boolean masking incurs in only a linear (in $x$) overhead whenever the performed 
operations on share-split polynomials are \texttt{xor} linear, nonlinear operations (e.g., Boolean \texttt{and}) have a quadratic overhead~\cite{ishai2003private}.

\begin{figure*}[!t]
\noindent\begin{minipage}{0.45\textwidth}
\footnotesize
    \begin{algorithm}[H]
        \caption{Masked HQC.PKE \textsf{Keygen}\label{alg:keygen}}
        \KwIn{none}
        \KwOut{$\kpub = (\mathbf{h},\mathbf{s})\in \mathcal{R}\times \mathcal{R}$\newline 
        $\kpri$$=$$(\mathbf{x},\mathbf{y})$$\in$$\mathcal{R}_\omega$$\times$$\mathcal{R}_\omega$}

        $\mathbf{h} \sample \mathcal{R}$\; 
        $(\mathbf{x},\mathbf{y}) \sample \mathcal{R}_\omega \times \mathcal{R}_\omega$\;
        $\ssplit{\mathbf{y}} \gets \text{\sc Split}(\mathbf{y})$\;
        $\ssplit{\mathbf{s}} \gets (\mathbf{h}\cdot\mathbf{y}_{,0},\ldots,\mathbf{h}\cdot\mathbf{y}_{,x-1})$\;
        $\ssplit{\mathbf{x}} \gets \text{\sc Split}(\mathbf{x})$\;
        $\ssplit{\mathbf{s}} \gets \ssplit{\mathbf{x}} \oplus \ssplit{\mathbf{s}}$  \;
        $\kpri \gets (\mathbf{x},\mathbf{y})$\;
        $\kpub \gets (\mathbf{h},\textsc{Combine}(\ssplit{\mathbf{s}}))$\;
        \KwRet $\kpri$, $\kpub$\;
        \BlankLine
        \BlankLine
        \BlankLine
        \BlankLine
        \BlankLine
        \BlankLine
        \BlankLine
        \BlankLine
        \BlankLine
        \BlankLine
        \BlankLine
    \end{algorithm}
\end{minipage}\hspace{0.5cm}
\begin{minipage}{0.5\textwidth}
\footnotesize
    \begin{algorithm}[H]
        \caption{Masked HQC.PKE \textsf{Encrypt}\label{alg:enc}}
        \KwData{$\mathbf{G}$: a $\mathds{F}_2^{k \times n}$ matrix generating an efficiently correctable code}
        \KwIn{$\kpub = (\mathbf{h},\mathbf{s})\in \mathcal{R}\times \mathcal{R}$\newline 
        $\mathbf{m}\in \mathbb{F}_2^k$}
        \KwOut{$(\mathbf{u}, \mathbf{v}) \in \mathcal{R} \times \mathcal{R}$}

        $\mathbf{e} \sample \mathcal{R}_{\omega_e}$, 
        $(\mathbf{r_1},\mathbf{r_2}) \sample \mathcal{R}_{\omega_r} \times \mathcal{R}_{\omega_r}$\;
        $\ssplit{\mathbf{r_1}} \gets \text{\sc Split}(\mathbf{r_1})$\;
        $\ssplit{\mathbf{r_2}} \gets \text{\sc Split}(\mathbf{r_2})$\;
        $\ssplit{\mathbf{u}} \gets \ssplit{\mathbf{r_1}} \oplus 
        (\mathbf{h}\cdot\mathbf{r_2}_{,0},\ldots,\mathbf{h}\cdot\mathbf{r_2}_{,x-1})$\;
        $\ssplit{\mathbf{m}} \gets \text{\sc Split}(\mathbf{m})$\;
        $\ssplit{\mathbf{v}} \gets \ssplit{\mathbf{m}}\mathbf{G}  \oplus \text{\sc Split}(\mathbf{e}) $\;
        $\ssplit{\mathbf{v}} \gets \ssplit{\mathbf{v}} \oplus (\mathbf{s}\cdot\mathbf{r_2}_{,0},\ldots,\mathbf{s}\cdot\mathbf{r_2}_{,x-1}) $\;
        $\mathbf{u} \gets \textsc{Combine}(\ssplit{\mathbf{u}})$\;
        $\mathbf{v} \gets \textsc{Combine}(\ssplit{\mathbf{v}})$\;
        \KwRet $(\mathbf{u},\mathbf{v})$\;
    \end{algorithm}
\end{minipage}
\vspace{-0.4cm}
\end{figure*}
To the end of devising our algorithmic proposal for efficient Boolean masking
of HQC.PKE, we observe that the following two properties hold for its structure:
\begin{prop}[Public operands]\label{prop:one} All the polynomial multiplications have a factor 
which is public and either input independent or constant (across a computation). 
It is thus possible to avoid the share-splitting of the public operand and 
rewrite the multiplication as a sequence of $x$ share-wise \texttt{xor} 
additions of the share-split operand and $x$ share-wise circular shifts.
\end{prop}
\begin{prop}[Share-split decoding]\label{prop:two} Computing in a share-split 
representation the decoding w.r.t. the public code $\mathbf{G}$ is  
performed in a more efficient way exploiting the linearity of the encoding: 
the $x-1$ random shares for the erroneous codeword to be decoded are obtained 
encoding $x-1$ random information words.
\end{prop}
\noindent\textbf{Masked} \textsf{Keygen}.
The masking of the key generation algorithm (Algorithm~\ref{alg:keygen}) is 
performed splitting the private key ($\kpri = (\mathbf{x},\mathbf{y}) \in 
\mathcal{R}_\omega \times \mathcal{R}_\omega$) element $\mathbf{v}$ (line $3$), 
and multiplying it by the public polynomial $\mathbf{h}$, which does not need
to be share-split (Property~\ref{prop:one}). To add the secret value 
$\mathbf{x}$ to $\ssplit{\mathbf{s}}$, we share-split it (line $5$) in turn 
and perform a share-wise \texttt{xor}. Finally, we combine the shares of 
$\ssplit{\mathbf{s}}$ to obtain the second half of the public key.

\noindent\textbf{Masked} \textsf{Encrypt}. In the encryption algorithm, 
we protect both the message $\mathbf{m}$ and the ephemeral values 
$\mathbf{r_1},\mathbf{r_2},\mathbf{e}$ by share splitting them (lines $2,3,5,6$).
The multiplications at lines $4$ and $6$ can be performed without share splitting
the public factors $\mathbf{h}$ and $\mathbf{s}$ -- by Property~\ref{prop:one}, 
indeed, they are the two components of the public key -- and therefore their 
share-wise multiplication only takes linear time in the number of shares $x$.
The encoding of the shared message $\ssplit{\mathbf{m}}$ is performed via a sequence
of linear operations, i.e. the share-wise multiplication by the public binary 
matrix $\mathbf{G}$, which therefore does not need to be share split, or any 
equivalent encoding technique.
\begin{figure*}[!t]
\noindent\begin{minipage}{0.99\textwidth}
\footnotesize
\begin{algorithm}[H]
    \caption{Masked HQC.PKE \textsf{Decrypt}\label{alg:dec}}
    \KwData{$\mathbf{G}$: a $\mathbb{F}_2^{k \times n}$ matrix generating an efficiently correctable code}
    \KwIn{$\kpri = (\mathbf{x},\mathbf{y}) \in \mathcal{R}_\omega \times \mathcal{R}_\omega$\newline $(\mathbf{u}, \mathbf{v}) \in \mathcal{R} \times \mathcal{R}$}
    \KwOut{$\ssplit{\mathbf{b}}$, such that $\text{\sc Combine}(\ssplit{\mathbf{b}})=\mathbf{m}\in \mathbb{F}_2^k$}

    $\ssplit{\mathbf{y}} \gets \text{\sc Split}(\mathbf{y})$\;
    $\ssplit{\mathbf{uy}} \gets (\mathbf{u}\cdot\mathbf{y}_0,\ldots,\mathbf{u}\cdot\mathbf{y}_{x-1})$\;
    $\mathbf{b} \sample \mathbb{F}_2^k$\;
    $\ssplit{\mathbf{b}} \sample \text{\sc Split}(\mathbf{b})$\;
    $\ssplit{\mathbf{c}} \gets \ssplit{\mathbf{b}}\mathbf{G}$\; 
    $\mathbf{a} \gets \text{\sc Collapse}(\mathbf{v} ,\ssplit{\mathbf{c}})$\;
    $\overline{\mathbf{b}} \gets \text{\sc Decode}(\text{\sc Collapse}(\mathbf{a},\ssplit{\mathbf{uy}}))$\tcp*[l]{\color{blue}$\overline{\mathbf{b}} = \mathbf{m} \oplus \mathbf{b} $}
    $\mathbf{b}_0 \gets \overline{\mathbf{b}} \oplus\mathbf{b}_0$\tcp*[l]{\color{blue}this sets $\ssplit{\mathbf{b}}=(\mathbf{m} \oplus \mathbf{b} \oplus\mathbf{b}_0,\mathbf{b}_1, \ldots,\mathbf{b}_{x-1} )$}
    \KwRet $\ssplit{\mathbf{b}}$\;
    \end{algorithm}
\end{minipage}
\vspace{-0.4cm}
\end{figure*}

\noindent\textbf{Masked} \textsf{Decrypt}. The decryption function needs to 
protect the private key component $\mathbf{y}$, which is split (line $1$ of 
Algorithm~\ref{alg:dec}), and subsequently multiplied by the public ciphertext
component $\mathbf{u}$, which does not need splitting (Property~\ref{prop:one}).
We now use Property~\ref{prop:two} to protect the decoding procedure (line $6$)
without computing $x$ times the decoding algorithm, which is the most demanding
portion of HQC.PKE \textsf{Decrypt}.
To this end, we denote with a helper procedure, \textsc{Collapse}$(\mathbf{a}, \ssplit{\mathbf{b}})$ the \texttt{xor} sum of $\mathbf{a}$ with all the shares of
$\mathbf{b}$, done in storage order of the shares themselves, as doing a hierarchical
reduction would lower the number of side channel samples to know for an attacker~\cite{ishai2003private,prouff2013masking}.
We protect the input of the \textsc{Decode} procedure, adding a random codeword
to it. This random codeword is obtained drawing a random information word 
$\mathbf{b}$ (line $4$), share splitting it (line $5$) and encoding all the shares, 
obtaining the share-split codeword $\ssplit{\mathbf{c}}$ (line $6$).
We note that a single \textsc{Decode} computation is performed, as opposed
to a straightforward masking scheme which would require $x$ of them to be 
computed.
Thanks to Property~\ref{prop:two}, decoding the sum of share-split codewords
will result in the sum of the corresponding share-split information words.
We therefore collapse the share-split codeword $\ssplit{\mathbf{c}}$ onto the
public decoder input, $\mathbf{v}$, while also performing the subtraction 
$\mathbf{v}-\mathbf{uy}$ (lines $6$,$7$, recalling that the subtraction 
among binary polynomials amounts to a coefficient-wise \texttt{xor}).
The result of the \textsc{Decode} procedure $\mathbf{b}_0$ is thus the message 
information word $\mathbf{m}$, with a random information word $\mathbf{b}$ 
added to it.
Adding such a value to the first share of the share-split representation of
$\ssplit{\mathbf{b}}= (\mathbf{b}_0,\mathbf{b}_1, \ldots,\mathbf{b}_{x-1} )$ turns such a representation into a share split representation of
$\mathbf{m}=(\mathbf{m} \oplus \mathbf{b} \oplus\mathbf{b}_0,\mathbf{b}_1, \ldots,\mathbf{b}_{x-1} )$ which is the value returned by the protected HQC.PKE \textsf{Decrypt}.
Such a value would be recombined, in case the HQC.PKE were used alone; instead, 
considering the HQC.KEM, can be forwarded to the protected components of the FO
construction for further processing.
